现在证明 6.1.2. 由 6.1.4 以及 (6.1.5) 的交换性, 只须证明\underline{对于} $D_{\bar\eta}$ 的所有局部单值同态 (6.1.4.1) 的第一个箭头,
\begin{equation}\tag{6.1.7.1}
\prod_{\ell\ne\operatorname{char}(k)}\mathbb Z_\ell(1)
\longrightarrow
\pi_1^{\mathrm{mod}}(D_{\bar\eta}-D_{\bar\eta}\cap Z,\bar a),
\end{equation}
在与下列同态复合后彼此共轭:
\begin{equation}\tag{6.1.7.2}
\pi_1^{\mathrm{mod}}(D_{\bar\eta}-D_{\bar\eta}\cap Z,\bar a)
\longrightarrow
\pi_1^{\mathrm{mod}}(D_\eta-D_\eta\cap Z,\bar a).
\end{equation}
为证明这一点, 注意对 $k(\eta)$ 的每个有限 Galois 扩张 $K$, 若
\[
k(\eta)\subset K\subset k(\bar\eta),
\]
则概形
\[
D_K\stackrel{\mathrm{def}}{=}D_\eta\times_{k(\eta)}K
\]
为 $D_\eta$ 的平展覆盖. 因而 $D_K-D_K\cap Z$ 为 $D_\eta-D_\eta\cap Z$ 的驯覆盖, 从而 $\pi_1^{\mathrm{mod}}(D_\eta-D_\eta\cap Z,\bar a)$ 在概形 $D_K\cap Z$ 上作用. 我们取 $K$ 充分大, 使 $D_K\cap Z$ 的每个点都为 $K$-有理点; 这是可行的, 因为由假设 $D_\eta\cap Z$ 在 $k(\eta)$ 上平展.

尚须证明 $\pi_1^{\mathrm{mod}}(D_\eta-D_\eta\cap Z,\bar a)$ 在 $D_K\cap Z$ 的点上的作用\underline{可迁}. 这将证明, 与 $D_K\cap Z$ 中不同点相伴的复合局部惯性同态
\[
\prod_{\ell\ne\operatorname{char}(k)}\mathbb Z_\ell(1)
\longrightarrow
\pi_1^{\mathrm{mod}}(D_\eta-D_\eta\cap Z,\bar a)
\]
彼此共轭. 而在 $D_K\cap Z$ 上的作用通过典范商分解:
\footnote{印本中 Galois 群的分子空缺. 紧接在公式之前的域塔决定它为 $k(\bar\eta)$.}
\[
\pi_1^{\mathrm{mod}}(D_\eta-D_\eta\cap Z,\bar a)
\longrightarrow
\pi_1(\eta,\bar\eta)
=\operatorname{Gal}\bigl(k(\bar\eta)/k(\eta)\bigr).
\]
$\operatorname{Gal}(k(\bar\eta)/k(\eta))$ 在 $D_K\cap Z$ 上的作用与包含映射
\[
D_K\cap Z\lhook\joinrel\longrightarrow\mathbb P^r_k\times_k K,
\]
相容, 其中 Galois 群在 $\mathbb P^r_k\times_k K$ 上仅通过第二个因子作用. 因而所需断言等价于说 $D_\eta\cap Z$ 仅有一个点, 即它为一个\underline{域}的谱. 用初等语言说, 这恰是下面的 Bertini 型事实 (取 $Z$ 的一个仿射方程 $F=0$, 并假设 $\bar a$ 为原点).

\subsubsection*{事实 6.1.8.}
设 $k$ 为代数闭域, $r\ge2$, 且 $F(X_1,\ldots,X_r)\in k[X_1,\ldots,X_r]$ 不可约, 并满足上文的条件 $F(0,\ldots,0)\ne0$. 则
\[
G(T)=F(X_1T,\ldots,X_rT)\in k(X_1,\ldots,X_r)[T]
\]
作为以 $k(X_1,\ldots,X_r)$ 为系数的 $T$ 的多项式不可约.

至此完成了 6.1.2 的证明.

\begin{editorialnote}
底本将上述一般直线的局部惯性同态误指为 (6.1.3); 该式实际是特化同态. 所指局部惯性箭头是 (6.1.4.1) 的第一个箭头, 亦即 (6.1.7.1), 此处订正引用. 事实 6.1.8 还须带上 $F(0,\ldots,0)\ne0$: 上文 $\bar a$ 为原点且 $\bar a\notin Z$ 已蕴含这一点. 若脱离此条件, 例如 $F=X_1+X_2^2$, 则 $F(X_1T,X_2T)=T(X_1+X_2^2T)$ 可约, 与所列事实不符. 本文在陈述中明确写出所继承的条件.
\end{editorialnote}

\subsubsection*{注 6.1.9.}
还有一种证明 6.1.2 的方法, 适用于一般得多的条件: 只须建立“单值同态”的一般相容性. 6.1.2 中的复合同态都共轭于 $Z$ 的泛点处的单值同态, 因而彼此共轭.

\subsubsection*{推论 6.1.10.}
设 $X$ 为代数闭域 $k$ 上的 $n$ 维光滑连通射影簇. 假设 $n=\dim X$ 为偶数, 或 $\operatorname{char}(k)\ne2$.\footnote{印本写作 $\operatorname{card}(k)\ne2$; (6.3.2) 中的对应假设以及驯单值论证要求的是特征.}
设 $D$ 为 $X$ 的一个 Lefschetz 线束, $\eta$ 为其泛点, $\bar\eta$ 为 $\eta$ 上的一个几何点. 则 $\pi_1(D-D\cap\check X,\bar\eta)$ 在 $H^{n-1}(X_{\bar\eta},\mathbb Q_\ell)$ 上的自然表示通过 $\pi_1^{\mathrm{mod}}(D-D\cap\check X,\bar\eta)$ 分解. 若此表示非平凡, 则 $D\cap\check X\ne\varnothing$, 且 $D\cap\check X$ 中每个点处的惯性群均非平凡地作用.

\subsubsection*{证明.}
若 $\dim(\check X)\le r-2$, 则 $D\cap\check X$ 为空, 由于 $\mathbb P^1$ 单连通, 没有什么要证明的. 否则, $\check X$ 为 $\mathbb P^r$ 中的超曲面 [XVII, 3.1.4], 由 3.2.10 及 XVII 3.5, 可取 $Z=\check X$、$D_o=D$ 应用 6.1.1 和 6.1.2. $\pi_1$ 在 $H^{n-1}(X_{\bar\eta},\mathbb Q_\ell)$ 上的作用为驯作用, 这一事实已由局部理论中的 Picard--Lefschetz 公式 [XV, 3.4] 建立, 下面将予以回顾. 最后, 非平凡的事实是, $\pi_1^{\mathrm{mod}}(D-D\cap\check X,\bar\eta)$ 由 $D\cap\check X$ 中各点处的惯性群生成 [SGA 1, XIII 2.10], 由此得到所需结论.

现在利用前面的结果研究 (5.3.5) 的条件 (A). 先回顾 Picard--Lefschetz 公式 [XV, 3.4] 的以下推论.
